% from ega2-1-fr.tex
\oldpage[II]{160}
de type fini, on sait (\hyperref[II.2.1.6-fr]{2.1.6}, (ii)) qu'il existe
$h>0$ et $m_0>0$ tels que $S_{n+h}=S_hS_n$ pour $n\geq m_0$~; on peut
évidemment supposer $m_0\geq h$. En outre, les $S_m$ sont des $S_0$-modules
de type fini (\hyperref[II.2.1.6-fr]{2.1.6}, (i)). Cela étant, si
$n\geq m_0+h$, on a $S_n^\natural=S_hS_{n-h}^\natural
=S_h^\natural S_{n-h}^\natural$~; et si $m<m_0+h$, on a, en posant
$E=S_{m_0}+\cdots+S_{m_0+h-1}$,
\[
  S_m^\natural=S_m+\cdots+S_{m_0+h-1}+S_hE+S_h^2E+\cdots.
\]
Pour $1\leq m\leq m_0$, soit $G_m$ la réunion de systèmes finis de
générateurs des $S_0$-modules $S_i$, pour $m\leq i\leq m_0+h-1$,
considérée comme partie de $S_{[m]}$. Pour $m_0+1\leq m\leq m_0+h-1$, soit
de même $G_m$ la réunion de systèmes finis de générateurs des $S_0$-modules
$S_i$, pour $m\leq i\leq m_0+h-1$, et de $S_hE$, considérée comme partie de
$S_{[m]}$. Il est clair que l'on a $S_m^\natural=S_0^\natural G_m$ pour
$1\leq m\leq m_0+h-1$, et par suite la réunion $G$ des $G_m$ pour
$1\leq m\leq m_0+h-1$ est un système de générateurs de la
$S_0^\natural$-algèbre $S^\natural$. On en conclut que si $S=S_0^\natural$
est un anneau noethérien, il en est de même de $S^\natural$.

\item[\rm(ii)] Il est clair que si $S_{n+1}=S_1S_n$ pour $n\geq n_0$, on a
$S_{n+1}^\natural=S_1S_n^\natural$, et \emph{a fortiori}
$S_{n+1}^\natural=S_1^\natural S_n^\natural$ pour $n\geq n_0$.
Inversement, cette dernière relation s'écrit
\[
  S_{n+1}+S_{n+2}+\cdots=(S_1+S_2+\cdots)(S_n+S_{n+1}+\cdots)
\]
et la comparaison des termes de degré $n+1$ (dans $S$) aux deux membres donne
$S_{n+1}=S_1S_n$.

\item[\rm(iii)] Si $S_n=S_1^n$ pour $n\geq n_0$, on a
$S_n^\natural=S_1^n+S_1^{n+1}+\cdots$~; comme $S_1^\natural$ contient
$S_1+S_1^2+\cdots$, on a $S_n^\natural\subset S_1^{\natural n}$, et par suite
$S_n^\natural=S_1^{\natural n}$ pour $n\geq n_0$. Inversement, les seuls
termes de $S_1^{\natural n}=(S_1+S_2+\cdots)^n$ qui soient de degré $n$ dans
$S$ sont ceux de $S_1^n$~; la relation $S_n^\natural=S_1^{\natural n}$
implique donc $S_n=S_1^n$.

\item[\rm(iv)] Il suffit de prouver que si un élément $g\in S_{k+h}$ est
considéré comme un élément de $S_k^\natural$ ($k>0$, $h\geq0$), alors il
existe un entier $n>0$ tel que dans $S_{kn}^\natural$, $g^n$ soit combinaison
linéaire des $f_\alpha^\natural$ à coefficients dans $S^\natural$. Par
hypothèse, il y a un entier $m_0$ tel que pour $m\geq m_0$, on ait, \emph{dans
$S$}, $g^m=\sum_\alpha c_{\alpha m}f_\alpha$, les indices $\alpha$ figurant
dans cette formule étant \emph{indépendants de $m$}~; en outre, on peut
évidemment supposer les $c_{\alpha m}$ homogènes, avec
\[
  \deg(c_{\alpha m})=m(k+h)-\deg f_\alpha
\]
dans $S$. Prenons alors $m_0$ assez grand pour que l'on ait
$km_0>\deg f_\alpha$ pour tous les $f_\alpha$ qui figurent dans $g^{m_0}$~;
pour tout $\alpha$, soit $c'_{\alpha m}$ l'élément $c_{\alpha m}$ considéré
comme de degré $km-\deg(f_\alpha)$ dans $S^\natural$~; on a alors, dans
$S^\natural$, $g^m=\sum_\alpha c'_{\alpha m}f_\alpha^\natural$, ce qui
termine la démonstration.
\end{enumerate}

\begin{env}[8.2.11]
\label{II.8.2.11-fr}
Considérons la $S_0$-algèbre graduée
\begin{equation}
\label{II.8.2.11.1-fr}
  S^\natural\otimes_S S_0=S^\natural/S_+S^\natural
  =\bigoplus_{n\geq0}S_{[n]}/S_+S_{[n]}.
\tag{8.2.11.1}
\end{equation}
Comme $S_n$ est un $S_0$-module quotient de $S_{[n]}/S_+S_{[n]}$, on a un
homomorphisme canonique de $S_0$-algèbres graduées
\begin{equation}
\label{II.8.2.11.2-fr}
  S^\natural\otimes_S S_0\to S
\tag{8.2.11.2}
\end{equation}
qui est évidemment \emph{surjectif}, et correspond par suite
(\hyperref[II.2.9.2-fr]{2.9.2}) à une \emph{immersion fermée} canonique
\begin{equation}
\label{II.8.2.11.3-fr}
  \operatorname{Proj}(S)\to
  \operatorname{Proj}(S^\natural\otimes_S S_0).
\tag{8.2.11.3}
\end{equation}
\end{env}

